% !TEX root = ../main.tex

% 预习报告	
\clearpage
\setcounter{section}{0}
\section{预习报告}
\vspace{-0.7cm} % 负值缩小间距
\noindent\textcolor{fgraygreen}{\rule{0.382\textwidth}{2pt} }
\vspace{7pt}
% 可用\ysySection{预习报告}一键实现

\subsection{实验概述}

通过测量电阻热噪声精密测量玻尔兹曼常量。

\subsection{实验原理}

实验利用锁相放大器测量噪声。时域上做等时间间隔采样，多次测量计算信号的\textbf{方均根值}。这句话的意思是：设备在等间隔的时间点上读取一个信号 $v(t_i) = s(t_i) + n(t_i)$，$s(t_i)$ 为 $t_i$ 时刻频率为 $f$ 的信号，$n(t_i)$ 为 $t_i$ 时刻的噪声，经与信号同频率 $f$ 的参考信号相敏检波后，输出 X 分量$X_f(t_i) = s_x(t_i) + n_f(t_i)$，是包含“同频”噪声的信号。定义噪声强度 X 分量测量值（这个式子显然是信号的方均根的表达式，同时它表达了噪声电压有效值）：
\[
v_{Nx}(f) = \sqrt{\frac{\sum_{i=1}^{M} \left(X_f(t_i) - \bar{X_f}\right)^2}{M}}
\]
其中，$\bar{X}_j = \frac{1}{M} \sum_{j=1}^{M} X_j$为输入信号的平均值，$M$ 为用于计算输入噪声值的样本数。

\begin{table}[h!]
	\centering
	\caption{单分量等效噪声带宽（ENBW）与陡降和时间常数$\tau$的对应关系}
	\label{tab:ENBW}
	\begin{tabular}{cccc}
		\hline
		\textbf{阶数(n)} & \textbf{Slope (dB/oct)} & \textbf{ENBW (1/$\tau$)} & \textbf{X 测量等待时间$t_X$}\\
		\hline
		1 & 6 & 1/4 & 4.6 \\
		2 & 12 & 1/8 & 6.6 \\
		3 & 18 & 3/32 & 8.4 \\
		4 & 24 & 5/64 & 10 \\
		\hline
	\end{tabular}
\end{table}

实际测量有各种各样的噪声，考虑按以下方式合成：
\[
v_{Nx}(f) = \sqrt{\frac{\sum_{i=1}^{M} \left( n_{f_x}(t_i) \right)^2}{M}} = \sqrt{\frac{\sum_{i=1}^{M} \left( \sum_{k=1}^{m} n_{f_xk}(t_i) \right)^2}{M}} 
\]

分离不同种类的噪声是困难的，除非在输入端做物理隔离。所以我们考虑的是非相干的噪声，按以下方式做合成：
\[
v_{Nx}(f) = \sqrt{\frac{\sum_{k=1}^{m} \sum_{i=1}^{M} \left( n_{f_xk}(t_i) \right)^2}{M}}
\]
上式给出了不同类型非相干噪声的叠加原理。

对随机噪声，X 与 Y 独立等价，即输入噪声功率为 X 或 Y 测量功率的 2 倍：
\[
v_N^2 = \frac{1}{M} \sum_{j=1}^{M} \left( n_{xT}^2(j) + n_{yT}^2(j) \right) = \frac{2}{M} \sum_{j=1}^{M} n_{xT}^2(j)
\]

\subsection{前思考题}
\begin{question}
	如何与同班组的同学分配频率点？
\end{question}
在分配频率点时，合理的分配能够确保实验的顺利进行，同时避免重复测量和资源浪费。首先，我们需要确定测量的频率范围和步进，通常根据实验要求和设备限制确定。如果讲义中没有明确规定频率范围，通常可以根据设备性能自行选择适合的范围，例如1 Hz到100 kHz。之后根据班组人数，合理分配频率点。假设一个实验组有4名同学，可以采用线性划分或对数划分的方式分配频率点。

在线性划分的方案中，频率点均匀分布在每个人负责的范围内。例如，在1 Hz到100 kHz的频率范围内测量20个频率点，每个同学负责其中的5个频率点。具体分配可以是：同学A测量1 Hz、10 Hz、100 Hz、1 kHz、10 kHz，同学B测量2 Hz、20 Hz、200 Hz、2 kHz、20 kHz，依此类推。这种分配方式简单易行，并确保每位同学都有从低频到高频的覆盖。

另一种方式是对数划分，适合频率跨度较大的实验。对于频率范围从1 Hz到100 kHz的情况，可以选择在对数坐标上均匀分配频率点。例如，16个频率点可以按对数划分，每个同学负责如1 Hz、10 Hz、100 Hz、1 kHz等倍数变化的频率点，这样的划分方式能够更好地覆盖高频段的变化。同时，低频点也会保持合理的间隔。

无论采用何种方案，都建议不同同学在某些频率点上进行交叉验证，以确保数据的一致性和准确性。例如，同学A和同学B可以在10 Hz和1 kHz进行重复测量，同学C和D则在50 Hz和5 kHz交叉测量。这样可以及时发现可能存在的误差，提高实验数据的可靠性。